\documentclass{situationist-preprint}

\title{\emph{No Exit} Reconsidered}
\msid{PHE-09-b}
\msfamily{PHE-09}
\msdiscipline{Phenomenology}
\mstype{Companion manuscript}
\mskeywords{Sartre; the look; No Exit; conflict; recognition; scarcity}

\begin{document}
\maketitle

\begin{abstract}
Sartre's \emph{No Exit} is usually read as a dramatization of the thesis that conflict is the original meaning of being-for-others, summed up in its most famous line, that hell is other people. PHE-09-a argues that the analysis of the look establishes objectification, not hostility, and that hostility arises only when objectification is degradation. We argue that the play supports this reading rather than the conflictual one. Its hell is built from four conditions that the play itself specifies: permanent exposure, confinement, stakes of self-worth placed in the other's verdict, and a scarcity of recognition among three. The look as such does none of the damage; the conditions do.
\end{abstract}

\section{The construction of the room}

In \emph{No Exit}, three dead people, Garcin, Inès and Estelle, are shown into a drawing room that is to be their hell \citep{Sartre1944}. There are no instruments of torture. The lights cannot be put out, the characters have no eyelids, the door is locked and, when it finally opens, none of them leaves. Each comes with a past that he or she needs to have judged favorably: Garcin, that he was not a coward; Estelle, that she is desirable and innocent; Inès, who claims to need nothing, needs Estelle. By the end Garcin draws the conclusion, that hell is other people.

On the conflictual reading of \emph{Being and Nothingness} \citep{Sartre1943}, the line states a general truth: since the look makes me an object and alienates my possibilities, every relation to the other is a struggle, and the play merely removes the distractions that usually hide it. PHE-09-a accepts the first step (Lemma~3.1, the look makes me an object for another), qualifies the second (Lemma~3.2, what the look takes is the monopoly of the meaning of my situation, not my possibilities), and denies the third: by Theorem~3.3, the look is hostile only if the objectification it effects is a degradation in the sense of Definition~2.3. The play is a test of this verdict, since it is designed to show the look at its worst.

\section{The four conditions}

We list the conditions that the play stipulates and ask which of them produce the hostility.

\begin{enumerate}
\item \emph{Exposure.} Each is always visible to the others. There is no sleep and no privacy.
\item \emph{Confinement.} The three cannot leave one another and cannot add anyone to the room.
\item \emph{Stakes.} Each has placed his or her worth in a verdict that only the others can give: Garcin's courage, Estelle's innocence and desirability.
\item \emph{Scarcity.} Each needs the recognition of one of the others, and the recognition each needs is withheld because it is being sought by, or given to, the third: Garcin needs Inès's verdict, Inès wants Estelle, Estelle wants Garcin.
\end{enumerate}

\begin{proposition}\label{p:conditions}
In \emph{No Exit}, the hostility of the look arises from conditions 3 and 4; conditions 1 and 2 intensify it but do not produce it.
\end{proposition}

\begin{proof}[Argument]
Consider the room with conditions 1 and 2 but without 3 and 4. Three people who are always visible to each other and confined together, but who have not staked their worth on a verdict and do not compete for scarce recognition, are objectified by one another at every moment. Nothing in the analysis of the look implies that this objectification ranks any of them below another. The play shows as much in the opening scenes, before the stakes are revealed, when the characters are embarrassed but not yet at war. The hostility begins when Inès's look takes on the character of a verdict on Garcin's cowardice, that is, when it ranks him below the standard on which he has staked his worth. It is then a degradation in the sense of Definition~2.3 of PHE-09-a. And it becomes inescapable through scarcity: each verdict that would release one of them is withheld because it is needed elsewhere in the triangle. Remove condition 3, and the look has nothing to degrade; remove condition 4, and a favorable verdict becomes available.
\end{proof}

Condition 4 is the notion of scarcity that Sartre later made central to his account of conflict in the \emph{Critique of Dialectical Reason} \citep{Sartre1960}, where antagonism among humans is tied to the material conditions under which there is not enough for all. The play anticipates that analysis in a single room: recognition is the scarce good, and the one from whom each needs it is turned towards someone else.

\section{Objections}

\begin{objection}
The play is meant to show the essence of being-for-others; its conditions are the essence made visible, not contingent additions.
\end{objection}

\begin{reply}
The conditions are not visible in the essence: they are stipulated by the dramatist, one by one, and the play depends on each. Sartre himself later said that the line had been misunderstood, and that it means that when relations with others are twisted or vitiated, the other can only be hell \citep{ContatRybalka1970}. Whatever weight one gives to an author's gloss, it agrees with what the construction of the room shows.
\end{reply}

\begin{objection}
Permanent exposure is already a degradation, since it removes the privacy without which one cannot have a standing of one's own.
\end{objection}

\begin{reply}
The loss of privacy is a harm, but it is a harm of social arrangement, not of the look: it consists in there being no time when one is not looked at. A single look does not remove privacy, and the harm of exposure grows with the arrangement, not with the objectification of any one look. That is what Theorem~3.3 of PHE-09-a predicts: hostility depends on conditions under which the look carries a ranking.
\end{reply}

\section{Conclusion}

\emph{No Exit} does not show that the look is essentially conflictual. It shows what the look becomes when people have staked their worth on one another's verdicts and when recognition is scarce. Under those conditions every look is a verdict and every verdict a degradation; without them, the look objectifies and nothing more, as PHE-09-a concludes.

\begin{thebibliography}{99}\small
\bibitem[Contat and Rybalka(1970)]{ContatRybalka1970} Contat, M. and Rybalka, M. (1970). \emph{Les écrits de Sartre}. Gallimard.
\bibitem[Sartre(1943)]{Sartre1943} Sartre, J.-P. (1943). \emph{L'être et le néant. Essai d'ontologie phénoménologique}. Gallimard.
\bibitem[Sartre(1944)]{Sartre1944} Sartre, J.-P. (1944). \emph{Huis clos}. Play in one act, first performed at the Théâtre du Vieux-Colombier, Paris.
\bibitem[Sartre(1960)]{Sartre1960} Sartre, J.-P. (1960). \emph{Critique de la raison dialectique}, vol.~1. Gallimard.
\end{thebibliography}

\end{document}
